\section{ĐÁNH GIÁ MÔ HÌNH}

\subsection{Các phương thức}
\begin{itemize}
    \item Hold-out (Chia train/test).
    \item Cross-validation (Kiểm thử chéo K-fold).
\end{itemize}

\subsection{Confusion Matrix (Ma trận nhầm lẫn)}
\textbf{Khái niệm:} Confusion Matrix là một bảng hai chiều dùng để đánh giá hiệu suất của mô hình phân loại. Nó hiển thị kết quả dự đoán của mô hình so với nhãn thực tế. Ma trận nhầm lẫn đặc biệt hữu ích trong việc phát hiện lỗi của mô hình và so sánh giữa các lớp với nhau.

\textbf{Các tham số:}
\begin{table}[H]
    \centering
    \renewcommand{\arraystretch}{1.3}
    \begin{tabular}{|l|c|c|}
        \hline
        & \textbf{Predicted Positive} & \textbf{Predicted Negative} \\
        \hline
        \textbf{Actual Positive} & True Positive (TP) & False Negative (FN) \\
        \hline
        \textbf{Actual Negative} & False Positive (FP) & True Negative (TN) \\
        \hline
    \end{tabular}
\end{table}

\begin{itemize}
    \item \textbf{True Positive (TP):} Số lượng các đối tượng thuộc lớp dương và được dự đoán đúng.
    \item \textbf{True Negative (TN):} Số lượng các đối tượng thuộc lớp âm và được dự đoán đúng.
    \item \textbf{False Positive (FP):} Số lượng các đối tượng thuộc lớp âm nhưng lại bị dự đoán là dương.
    \item \textbf{False Negative (FN):} Số lượng các đối tượng thuộc lớp dương nhưng lại bị dự đoán là âm.
\end{itemize}

\textbf{Ví dụ minh họa:} Xét 10 mẫu với nhãn thực tế và kết quả dự đoán của mô hình như sau:
\begin{table}[H]
    \centering
    \begin{tabular}{|c|c|c|}
        \hline
        \textbf{Mẫu} & \textbf{Nhãn thực tế} & \textbf{Dự đoán của mô hình} \\
        \hline
        1 & + & + \\ \hline
        2 & + & - \\ \hline
        3 & - & - \\ \hline
        4 & - & + \\ \hline
        5 & + & + \\ \hline
        6 & + & + \\ \hline
        7 & - & - \\ \hline
        8 & - & - \\ \hline
        9 & + & - \\ \hline
        10 & - & + \\ \hline
    \end{tabular}
\end{table}

Chọn lớp ``+'' là lớp dương. Đối chiếu nhãn thực tế với dự đoán, ta có:
\begin{itemize}
    \item $\text{TP (True Positive): Mẫu 1, 5, 6} \Rightarrow \text{TP} = 3$.
    \item $\text{TN (True Negative): Mẫu 3, 7, 8} \Rightarrow \text{TN} = 3$.
    \item $\text{FP (False Positive): Mẫu 4, 10} \Rightarrow \text{FP} = 2$.
    \item $\text{FN (False Negative): Mẫu 2, 9} \Rightarrow \text{FN} = 2$.
\end{itemize}

\textbf{Bước 4 -- Điền các giá trị vào ma trận nhầm lẫn:}
\begin{table}[H]
    \centering
    \renewcommand{\arraystretch}{1.3}
    \begin{tabular}{|l|c|c|}
        \hline
        & \textbf{Dự đoán dương (+)} & \textbf{Dự đoán âm (-)} \\
        \hline
        \textbf{Thực tế dương (+)} & 3 (TP) & 2 (FN) \\
        \hline
        \textbf{Thực tế âm (-)} & 2 (FP) & 3 (TN) \\
        \hline
    \end{tabular}
\end{table}

\textbf{Bước 5 -- Tính các chỉ số đánh giá:}
\begin{align*}
\text{Accuracy}
&= \frac{TP+TN}{TP+TN+FP+FN}
= \frac{3+3}{3+3+2+2}
= \frac{6}{10}
= 60\%,\\
\text{Precision}
&= \frac{TP}{TP+FP}
= \frac{3}{3+2}
= \frac{3}{5}
= 60\%,\\
\text{Recall}
&= \frac{TP}{TP+FN}
= \frac{3}{3+2}
= \frac{3}{5}
= 60\%,\\
\text{F1-Score}
&= \frac{2 \cdot \text{Precision} \cdot \text{Recall}}
{\text{Precision}+\text{Recall}}
= \frac{2 \cdot 0.6 \cdot 0.6}{0.6+0.6}
= 60\%.
\end{align*}
